% Einheit 4 von 4 – Sachgeraden (bank/geraden/e4.jsonl, zusammenbau v0.1)
%% TODO Zeitmarke Einheit 4: Marken-Zeile nicht lesbar
%% TODO Zweigzeile Teil 1: was hier gelernt wird (Satz in Schülersprache aus dem Einheitstitel) – Einheit 4
\einheitenkopf[e4]{Einheit 4 von 4 · Sachgeraden}
\zweigzeile{Abitur GK}
\setcounter{aufgabe}{17}

%% TODO Titel: Ich-kann-Satz fehlt in der Bank (Platzhalter: Kettenname) – Nr. 18
%% TODO Anweisung: ein Satz über der ersten Teilaufgabe fehlt in der Bank – Nr. 18
\begin{aufgabe}{Sachgeraden}
%% TODO \rechenplatz{n}: Schrittzahl des Grundfalls fehlt in der Bank (nur bei fester Schreibform, 2.3 b)
\begin{teile}
\teil Ein Seil ist zwischen zwei Masten gespannt. Wird es im Modell durch eine Gerade oder durch eine Strecke beschrieben? \kreuz{Gerade} \\ \kreuz{Strecke}
\teil Ein Seil ist an zwei Masten in $P(0 | 0 | 8)$ und $Q(6 | 8 | 5)$ befestigt (1 LE = 1 m). Gib eine Gleichung der Seilstrecke an.
\teil Ein unteres Seil ist geradlinig von $A(0 | 0 | 6)$ nach $B(4 | 8 | 4)$ gespannt, ein oberes verläuft entlang $h\colon \vec x = (-2 | 2 | 10) + t \cdot (3 | 0 | 0)$ (1 LE = 1 m). Ein Punkt des oberen Seils liegt senkrecht über einem Punkt des unteren. Berechne den Abstand dieser beiden Punkte. \hfill (Abitur 2022 GK) \\ \leerfeld[m]
\teil Ein Seil ist geradlinig zwischen zwei Masten gespannt; die Fußpunkte sind $F_1(0 | 0 | 0)$ und $F_2(6 | 8 | 0)$, befestigt ist es in 9 m und in $7{,}5$ m Höhe (1 LE = 1 m). Berechne die Neigung des Seils in Prozent. \hfill (Abitur 2022 GK) \\ \leerfeld \%
\teil Ein Sonnensegel wird mit einem geradlinig gespannten Seil befestigt. Das Seil beginnt am ersten Mast im Punkt $(0 | 0 | 3)$ und endet am zweiten Mast, der senkrecht durch $M(6 | 8 | 0)$ steht. Unterwegs liegt es auf der Kante $UV$ eines waagerechten Dachs in 4 m Höhe auf, mit $U(4 | 3 | 4)$ und $V(1 | 6 | 4)$; das Verhältnis, in dem der Auflagepunkt die Kante teilt, ist bekannt (1 LE = 1 m). Beschreibe ein Verfahren, mit dem sich der Abstand des Befestigungspunkts am zweiten Mast vom Dach berechnen lässt. \hfill (Abitur 2018 GK)
\end{teile}
\end{aufgabe}

%% TODO Titel: Ich-kann-Satz fehlt in der Bank (Platzhalter: Kettenname) – Nr. 19
%% TODO Anweisung: ein Satz über der ersten Teilaufgabe fehlt in der Bank – Nr. 19
\begin{aufgabe}{Horizontalabstand aus vorgegebener Neigung in Prozent und Höhendifferenz berechnen}
\begin{teile}
\teil Ein Steg beginnt am rechten Ufer in $0{,}9$ m Höhe über dem Wasserspiegel und endet links auf der Uferzone in $0{,}3$ m Höhe; er soll eine Steigung von 4\,\% haben. Im Querschnitt liegt der Anfang bei $x = 0$, die linke Uferzone beginnt bei $x = -10$ (Angaben in m). Wie weit liegt das linke Ende des Stegs vom Rand der Uferzone entfernt? \hfill (Abitur 2017 GK) \\ \leerfeld[m]
\end{teile}
\end{aufgabe}

%% TODO Titel: Ich-kann-Satz fehlt in der Bank (Platzhalter: Kettenname) – Nr. 20
%% TODO Anweisung: ein Satz über der ersten Teilaufgabe fehlt in der Bank – Nr. 20
\begin{aufgabe}{Sachgeraden – Fehler finden, Begründen, Anwendung}
\begin{teile}
\teil Finde den Fehler. Zwei Seile verlaufen entlang der windschiefen Geraden $g$ und $h$; gefragt ist der Abstand eines Punkts des oberen Seils vom senkrecht darunterliegenden Punkt des unteren Seils. Jana berechnet dafür den Abstand der windschiefen Geraden $g$ und $h$.
\teil Begründe, warum die Bedingung „der Punkt $Q$ liegt senkrecht über dem Punkt $P$“ zwei Gleichungen liefert.
\teil Ein Flugzeug steigt nach dem Start geradlinig: um 12.00 Uhr ist es in $A(0 | 0 | 0{,}5)$, eine Minute später in $B(4 | 3 | 1{,}1)$ (1 LE = 1 km). Nach wie vielen Minuten erreicht es die Reiseflughöhe von 5 km, wenn es gleichmäßig so weitersteigt? \leerfeld[min]
\end{teile}
\end{aufgabe}
